\scanpage{194}
ALGÈBRES de LIE SEMI-SIMPLES \mspage{(8)}\hfill Tribu n\textsuperscript{o} 59

Dire `` ces éléments $N_{\alpha\beta}$ '' à la p. 46, \unsure{l|1}.6 bas.

\uline{n\textsuperscript{o} 6} - Enoncer ainsi la prop. 2 : a) tout sous-espace $V$ stable par $\underline{h}$ est de la forme, (sous-espace $\underline{h}'$ de $\underline{h}$) $\oplus\sum_{\alpha\in\Sigma}\underline{g}^{\alpha}$ ; b) pour que $V$ soit une sous-algèbre, il faut et il suffit que $\Sigma$ soit clos et que $\alpha\in\Sigma\cap(-\Sigma)$ implique $H_\alpha\in\underline{h}'$. Exhiber alors le radical de $V$, (une partie vient de $\underline{h}$, une autre des $\alpha\in\Sigma$ tels que $-\alpha\notin\Sigma$).

Enoncé de la prop. 3 : $V$ réductive $\Longleftrightarrow\Sigma$ clos et symétrique ; $V$ semi-simple $\Longleftrightarrow\Sigma$\unsure{,|.} clos et symétrique et $\underline{h}'$ engendré par les $H_\alpha$ pour $\alpha\in\Sigma$ ; son système de racines est alors celui que dit la orop. 3.

Le cor. 2, p. \unsure{47|49}, ira en exer6. La dém. de la prop. 4 se simplifie, maintenant qu'on connait le radical. Dans le cor., p. 51, dire `` racines positives par rapport à une base de $R$ ''.

A propos des Borels, (p. 51 bas), dire que la forme quadratique est nulle sur $\underline{n}^{+}$ et $\underline{n}^{-}$, et que leurs orthogonaux sont $\underline{b}^{-}$ et $\underline{b}^{+}$. On parlera des paraboliques, (ou paraboréliques).

\uline{n\textsuperscript{o} 7} - Dans la prop. 7, indexer $S$, (qui s'appelle maintenant `` base de $R$ ''). A la l.8 bas, (p. 52), dire que les autres formules ``ont déjà été démontrées'', et donner les références.

Donner ici les relations $(\operatorname{ad}X_i)^{-a_{ij}+1}X_j=(\operatorname{ad}Y_i)^{-a_{ij}+1}Y_j=0$, (cf. \S\,4).

Pas de nom pour $(H_i,X_i,Y_i)$ en haut de la p. 53, ça conflicte ; on dira que ce système est associé à l'épinglage.

\uline{n\textsuperscript{o} 8} - Au début préciser que $X_\alpha\in\underline{g}^{\alpha}$. Le lemme 4 remonte au \S\,2, n\textsuperscript{o} 2. Introduire $p$ aussi dans l'énoncé de la prop. 6. La formule du bas de la p. 54 tient à ce que $\alpha+\beta$ est une racine ; le dire.
\margin{193}
